% ======================================================================
%  PHỤ LỤC
% ======================================================================
\chapter{Hướng dẫn tái lập thực nghiệm}
\label{ch:phuluc-repo}

Toàn bộ mã nguồn của đồ án được công bố tại kho lưu trữ:
\begin{center}
\url{https://github.com/ThinhSama234/HCMUS-data-mining-AutoML}
\end{center}

Kho mã nguồn kèm tệp \texttt{README.md} mô tả chi tiết cách thiết lập môi trường
và chạy thực nghiệm. Quy trình tái lập gồm các bước sau.

\paragraph{Bước 1 --- Chuẩn bị dữ liệu.} Tải các tệp thô từ Kaggle theo danh mục
đăng ký trong \texttt{scripts/prep.py}, đặt vào thư mục \texttt{data/raw/} rồi
chạy lệnh tiền xử lý. Kết quả là mỗi bộ dữ liệu được chuẩn hoá thành một tệp dữ
liệu và một tệp siêu dữ liệu mô tả cột nhãn, dạng tác vụ và độ đo:
\begin{center}
\texttt{python scripts/prep.py -{}-all}
\end{center}

\paragraph{Bước 2 --- Sinh phân hoạch fold.} Phân hoạch được sinh một lần và dùng
chung cho mọi thư viện, bảo đảm điều kiện công bằng theo
Định nghĩa~\ref{def:congbang}:
\begin{center}
\texttt{python scripts/gen\_folds.py -{}-n-splits 5 -{}-seed 42}
\end{center}

\paragraph{Bước 3 --- Thiết lập môi trường ảo cho từng thư viện.} Mỗi thư viện
AutoML được cài đặt vào một môi trường ảo riêng dưới thư mục \texttt{envs/venvs/}
nhằm loại trừ xung đột phụ thuộc giữa các thư viện.

\paragraph{Bước 4 --- Chạy thực nghiệm.} Tiến trình điều phối sinh và phân phối
toàn bộ công việc; tham số \texttt{-{}-time-budget} quy định ngân sách thời gian
cho mỗi lần huấn luyện:
\begin{center}
\texttt{python scripts/orchestrator.py -{}-time-budget 300}
\end{center}
Quá trình có thể được tiếp tục sau gián đoạn bằng cách truyền lại định danh lượt
chạy qua tham số \texttt{-{}-run-id}, hoặc giới hạn phạm vi bằng
\texttt{-{}-frameworks} và \texttt{-{}-datasets} khi chia việc cho nhiều máy.

\paragraph{Bước 5 --- Hợp nhất và phân tích.} Các tệp kết quả từng phần thu được
từ nhiều máy được gộp thành một báo cáo duy nhất, sau đó sinh bảng tổng hợp và
biểu đồ:
\begin{center}
\texttt{python scripts/merge.py reports/partial\_*.json}
\end{center}

\paragraph{Cấu hình tái lập.} Các tham số quyết định tính tái lập của kết quả
trình bày trong báo cáo được liệt kê ở Bảng~\ref{tab:cauhinh}.

\begin{table}[H]
  \centering
  \small
  \caption{Cấu hình thực nghiệm dùng để tái lập kết quả.}
  \label{tab:cauhinh}
  \renewcommand{\arraystretch}{1.2}
  \begin{tabularx}{\linewidth}{@{}l X@{}}
    \toprule
    \textbf{Tham số} & \textbf{Giá trị} \\
    \midrule
    Số fold $K$              & 5 \\
    Hạt giống ngẫu nhiên     & 42 \\
    Phép chia (phân loại)    & Phân tầng theo nhãn (\textit{stratified $K$-fold}) \\
    Phép chia (hồi quy)      & $K$-fold ngẫu nhiên \\
    Ngân sách thời gian      & 60 giây và 300 giây cho mỗi lần huấn luyện \\
    Thư viện so sánh         & AutoGluon, FLAML, H2O AutoML, MLJAR \\
    Số bộ dữ liệu            & 12 (5 nhị phân, 3 đa lớp, 4 hồi quy) \\
    \bottomrule
  \end{tabularx}
\end{table}

\chapter{Ghi chú về việc sử dụng công cụ hỗ trợ}
\label{ch:phuluc-ai}

Theo quy định của đồ án, nhóm ghi chú minh bạch những phần có sử dụng công cụ trí
tuệ nhân tạo trong quá trình thực hiện. Nhóm đã kiểm tra, hiệu chỉnh và chịu trách
nhiệm hoàn toàn về tính chính xác của mọi nội dung trình bày trong báo cáo cũng
như mã nguồn nộp kèm.

\begin{table}[H]
  \centering
  \small
  \caption{Ghi chú sử dụng công cụ hỗ trợ trong đồ án.}
  \label{tab:ai-usage}
  \renewcommand{\arraystretch}{1.25}
  \begin{tabularx}{\linewidth}{@{}l l X@{}}
    \toprule
    \textbf{Hạng mục} & \textbf{Công cụ} & \textbf{Mức độ sử dụng và cách kiểm chứng} \\
    \midrule
    {[}Soạn thảo báo cáo{]} & {[}tên công cụ{]} & {[}Mô tả phạm vi hỗ trợ; nhóm đã đọc, hiệu chỉnh và đối chiếu mọi số liệu với kết quả thực nghiệm gốc.{]} \\
    {[}Hỗ trợ lập trình{]}  & {[}tên công cụ{]} & {[}Mô tả phạm vi hỗ trợ; nhóm đã kiểm thử và xác nhận tính đúng đắn của mã nguồn.{]} \\
    {[}Trực quan hoá dữ liệu{]} & {[}tên công cụ{]} & {[}Mô tả phạm vi hỗ trợ; nhóm đã kiểm tra tính nhất quán giữa biểu đồ và dữ liệu gốc.{]} \\
    \bottomrule
  \end{tabularx}
\end{table}
